\documentclass[12pt]{article}
\usepackage{tcolorbox}
\usepackage{amsmath}


\include{preamble}
%All credit goes towards Dr. Adam Kapelner for providing the preamble and homework template.

\newtoggle{professormode}




\title{MATH 241 Spring 2026 Homework \#10 }

\author{Steven Ou} %STUDENTS: write your name here

\iftoggle{professormode}{
\date{Due via Brightspace by 11:59PM Tuesday, May 19, 2025 \\ \vspace{0.5cm} \small (this document last updated \today ~at \currenttime)}
}

\renewcommand{\abstractname}{Instructions and Philosophy}

\begin{document}
\maketitle

\iftoggle{professormode}{
\begin{abstract}
The path to success in this class is to do many problems. Unlike other courses, exclusively doing reading(s) will not help. Coming to lecture is akin to watching workout videos; thinking about and solving problems on your own is the actual ``working out.''  Feel free to \qu{work out} with others; \textbf{I want you to work on this in groups.}

The problems below are color coded: \ingreen{green} problems are considered \textit{easy} and marked \qu{[easy]}; \inorange{yellow} problems are considered \textit{intermediate} and marked \qu{[harder]}, \inred{red} problems are considered \textit{difficult} and marked \qu{[difficult]} and \inpurple{purple} problems are extra credit. The \textit{easy} problems are intended to be ``giveaways'' if you went to class. Do as much as you can of the others; I expect you to at least attempt the \textit{difficult} problems. 

Bonus points are given as a bonus if the homework is typed using \LaTeX. Links to installing \LaTeX~and programs for compiling \LaTeX~is found on the syllabus. You are encouraged to use \url{https://www.overleaf.com/}. If you are handing in homework this way, read the comments in the code; there is a line to comment out and you should replace my name with yours. The easiest way to use Overleaf is to go to \textit{File} and select \textit{Make a Copy}. If you are asked to make drawings, you can take a picture of your handwritten drawing and insert them as figures or leave space using the \qu{$\backslash$vspace} command and draw them in after printing or attach them stapled.


\end{abstract}

\thispagestyle{empty}
\vspace{1cm}
NAME: \line(1,0){380}
\clearpage
}

% \problem{In this problem we will review the Poisson random variable covered in Section 4.6.}

% \begin{enumerate}

% \inthenotessubproblem{What is the pmf of the Poisson random variable?}

% \vspace{25mm}


% \inthenotessubproblem{What is the support of the Poisson random variable?}

% \vspace{10mm}


% \inthenotessubproblem{Under what conditions do we obtain a Poisson random variable? We discussed two ways in class -- just list the simpler way. }

% \vspace{25mm}

% \hardsubproblem{The number of times that a person contracts a cold in a given year is a Poisson random variable with parameter $\lambda = 5$. Suppose that a new wonder drug (based on large quantities of vitamin C) has just been marketed that reduces the Poisson parameter to  $\lambda = 3$ for 75\%  percent of the population. For the other 25\%  percent of the population, the drug has no appreciable effect on colds. If an individual tries the drug for a year and has 2 colds in that time, how likely is it that the drug is beneficial for him or her?}

% \newpage


% \easysubproblem{(Optional): What is a whitepaper? Feel free to search this term on the Internet.}

% \vspace{20mm}

% \easysubproblem{(Optional): Who is Satoshi Nakamoto?  }

% \vspace{20mm}

% \easysubproblem{(Optional): True or False? In the Bitcoin whitepaper (published in 2008), the Poisson distribution is used.}

% \end{enumerate}

\problem{In this problem, we will review the ideas discussed in Section 5.1.}
\begin{enumerate}
\inthenotessubproblem{What are some quantities that cannot be modeled by a discrete random variable?}

\vspace{20mm}

\easysubproblem{True or False? Given a random variable $X$, it is possible that $X$ does not have a probability mass function.}

\inthenotessubproblem{Define what is meant by a \textit{continuous} random variable.}

\newpage

\intermediatesubproblem{ Suppose that $X$ is a continuous random variable whose probability density function is given by 
$$
f(x) = 
\begin{cases}
c(4x - 2x^2) & \text{if} ~ 0< x  < 2  \\   
0 & \text{otherwise} \\
\end{cases} 
$$
What is the value of $c$? Hint: Apply Humpty-Dumpty.}

\vspace{30mm}

\intermediatesubproblem{ Suppose that $X$ is a continuous random variable whose probability density function is given by 
$$
f(x) = 
\begin{cases}
c(4x - 2x^2) & \text{if} ~ 0< x  < 2  \\   
0 & \text{otherwise} \\
\end{cases} 
$$
Find the cumulative distribution function of $X$.}

\vspace{90mm}

\intermediatesubproblem{ Suppose that $X$ is a continuous random variable whose probability density function is given by 
$$
f(x) = 
\begin{cases}
c(4x - 2x^2) & \text{if} ~ 0< x  < 2  \\   
0 & \text{otherwise} \\
\end{cases} 
$$
Find $\mathbb{E}[X]$.}

\vspace{30mm}

\intermediatesubproblem{ Suppose that $X$ is a continuous random variable whose probability density function is given by 
$$
f(x) = 
\begin{cases}
c(4x - 2x^2) & \text{if} ~ 0< x  < 2  \\   
0 & \text{otherwise} \\
\end{cases} 
$$
Find $\mathbb{E}[X^2]$.}

\vspace{30mm}

\hardsubproblem{Suppose that a random variable $X$ measures the lifetime of a certain type of electronic device in hours. Suppose its probability density function is given by:
 $$
f(x) =
\begin{cases}
\frac{10}{x^2}, & \text{if} ~ x>10  \ \\ 
0 , & \text{otherwise} \\
\end{cases}
$$ 
What is the probability that of 6 such types of devices, at least 3 will function for at least 15 hours? (You may assume that the devices operate independently of one another). You may leave your answer in terms of a calculator command. }

\end{enumerate}

\newpage

\problem{These exercises will review the king of all random variables -- the Normal random variable.}

\begin{enumerate}

\inthenotessubproblem{What is the density of the normal random variable?}

\vspace{10mm}

\inthenotessubproblem{What is the support of the normal random variable?}

\vspace{5mm}

\easysubproblem{What is the cumulative distribution function of the normal random variable?}
\vspace{10mm}

\easysubproblem{If $X \sim \mathcal{N}(\mu, \sigma^2)$ then what is $\mathbb{E}[X]$?}

\vspace{5mm}

\easysubproblem{If $X \sim \mathcal{N}(\mu, \sigma^2)$ then what is $\mathbb{V}ar[X]$?}

\vspace{5mm}

\easysubproblem{What is the standard normal random variable?}

\vspace{5mm}

\inthenotessubproblem{What is the density of the standard normal random variable? Use the appropriate letter/notation.}

\vspace{10mm}

\inthenotessubproblem{What is the cumulative distribution function of the standard normal random variable? Use the appropriate letter/notation.}

\vspace{10mm}

\easysubproblem{Suppose $Z \sim \mathcal{N}(0,1)$. Find $P(Z<2.41)$. You may leave your answer in terms of an integral, in terms of $\Phi$, in terms of a calculator command, or in terms of an actual number rounded to four decimal places. } 

\vspace{25mm}

\easysubproblem{Suppose $Z \sim \mathcal{N}(0,1)$. Find $P(Z> 2.41)$. You may leave your answer in terms of an integral, in terms of $\Phi$, in terms of a calculator command, or in terms of an actual number rounded to four decimal places.}

\vspace{15mm}


\easysubproblem{Suppose $Z \sim \mathcal{N}(0,1)$. Find $P(Z = 2.41)$.}

\vspace{5mm}


\easysubproblem{Suppose $X \sim \mathcal{N}(241,17)$. Find $P(X > 231)$. You may leave your answer in terms of an integral, in terms of $\Phi$, in terms of a calculator command, or in terms of an actual number rounded to four decimal places.}

\vspace{15mm}


\easysubproblem{Suppose $X \sim \mathcal{N}(241,17)$. Find $P(X < 231)$. You may leave your answer in terms of an integral, in terms of $\Phi$, in terms of a calculator command, or in terms of an actual number rounded to four decimal places.}
\end{enumerate}

\newpage 

\problem{These exercises are designed to review the first half of Section 5.4. Suppose the joint  mass function of $X$ and $Y$ is given by: \begin{figure}[h!]
    \centering
    \includegraphics[width=0.5\linewidth]{joint_distribution.png}
    \label{fig:enter-label}
\end{figure}}

\begin{enumerate}

\easysubproblem{Define what is meant by a joint mass function (jmf). }

\vspace{20mm}

\intermediatesubproblem{Find $P(X \leq 2, Y \leq 2)$.}

\vspace{20mm}

\intermediatesubproblem{Find $f_X(2)$.}

\vspace{20mm}

\intermediatesubproblem{Find $f_Y(4)$.}

\vspace{20mm}


\intermediatesubproblem{Find $f_X(x)$.}

\newpage

\intermediatesubproblem{Find $f_Y(y)$.}

\vspace{40mm}

\inthenotessubproblem{What does it mean for two random variables, $X$ and $Y$, to be independent?}

\vspace{40mm}

\hardsubproblem{For the joint probability mass function specified above, is $X$ and $Y$ independent? Explain why or why not.}

\vspace{40mm}

\end{enumerate}

\problem{In this problem, we will review the last part of Section 5.4.}

\begin{enumerate}

\easysubproblem{Explain what is meant by a convolution of random variables.}

\vspace{25mm}

\intermediatesubproblem{Suppose $X_1, X_2 \iid \geometric{p}$}. Let $T = X_1 + X_2$. What is the support of $T$?

\vspace{10mm}

\hardsubproblem{Find $P(T = 2)$.}

\vspace{25mm}

\hardsubproblem{Find $P(T = 3)$.}

\vspace{30mm}

\hardsubproblem{Find $P(T = 4)$.}

\vspace{35mm}

\hardsubproblem{Based off the last three computations, can you guess the probability mass function of $T$? You do not have to mathematically prove your answer. (If you are curious, the actual proof requires induction and the general convolution formula which you are not required to know).}

\vspace{25mm}

\end{enumerate}


\problem{Here is the last problem for the semester. Congratulations on making it this far! These exercises are designed to review Section 5.5.}
\begin{enumerate}

\easysubproblem{Without simply quoting the formal statement of the Central Limit Theorem, can you explain what the Central Limit Theorem states?}

\newpage 

\hardsubproblem{ Suppose the random variables, $X_1, X_2, \ldots, X_{100}$ are independent and all have the following density:
\begin{Large}
$$
f(x) =
\begin{cases}
\frac{x^2}{72} & \text{if} ~ 0<x<6 \ \\ 
0 , & \text{otherwise} \\
\end{cases}
$$ 
\end{Large}

If $Y = X_1 + X_2 + \ldots + X_{100}$, then find/approximate $P(439 \leq Y \leq 461)$. You may leave your answer in terms of an integral, in terms of $\Phi$, in terms of a calculator command, or in terms of an actual number rounded to four decimal places.
 }

 \vspace{70mm}

\hardsubproblem{Suppose the random variables $X_1, X_2, \ldots, X_{36}$  are independent and all  have the density specified below. 

\begin{Large}
\[
  f(x) =
  \begin{cases}
    \frac{1}{20}x^3 & \text{if $1<x<3$} \\
    0 & \text{otherwise} 
  \end{cases}
\]
\end{Large}

Find/approximate  $P(2 \leq \overline{X}_{36} \leq 2.5)$.   You may leave your answer in terms of an integral, in terms of $\Phi$, in terms of a calculator command, or in terms of an actual number rounded to four decimal places.
}

\end{enumerate}

\end{document}
